% GENERATED FILE. Edit canonical lessons, then run npm run generate:books.
% canonical-source-hash: fnv1a64:53e77327405f628c
% canonical-lessons: SA-C484-ayurvedah, SA-C484-vatah, SA-C484-pittam, SA-C484-kaphah, SA-C484-pathyam

\chapter{Ayurveda}
\label{ch:sa-ayurveda}
\hlchaptermodality{484}

\begin{chapteropening}
\textbf{By the end of this chapter:} \emph{I can name the three humours of Ayurveda.}

Complete the last lesson of chapter 484: I can name the three humours of Ayurveda.
\end{chapteropening}

\section[\texorpdfstring{\sk{आयुर्वेदः} (āyurvedaḥ)}{आयुर्वेदः (āyurvedaḥ)}]{\sk{आयुर्वेदः} (āyurvedaḥ) — the science of health; Ayurveda}

\label{lesson:SA-C484-ayurvedah}



\begin{quote}
Before the new one: say the Sanskrit for a decoction, then the Sanskrit for smears.
\end{quote}

\subsection*{You'll want to know: \sk{आयुर्वेदः}}
\textbf{\sk{आयुर्वेदः}} — \emph{āyurvedaḥ} — \textquotedblleft{}the science of health; Ayurveda\textquotedblright{}.

\begin{quote}
\sk{आयुर्वेदः} \sk{आरोग्यस्य} \sk{ज्ञानम्} \sk{अस्ति।}
\end{quote}

\begin{quote}
\emph{āyurvedaḥ ārogyasya jñānam asti.}
\end{quote}

\begin{quote}
Ayurveda is the knowledge of health.
\end{quote}

A masculine noun.

\begin{cousinweb}
\textbf{\sk{आयुर्वेदः}} (\emph{āyurvedaḥ}) is built from \emph{āyus}, \textquotedblleft{}life\textquotedblright{}, and \emph{veda}, \textquotedblleft{}knowledge\textquotedblright{}, from the root √vid; MW gives \textquotedblleft{}the science of health or medicine\textquotedblright{} (MW).
\end{cousinweb}

\begin{grammarlens}[title={What you've built}]
One more word for describing what you lived and felt.
\end{grammarlens}

\subsection*{Your turn}
\begin{itemize}
\raggedright
  \item \emph{Say it:} \emph{āyurvedaḥ}
  \item \emph{Say it:} \emph{āyurvedaḥ}, once more
  \item \emph{Say it:} \emph{āyurvedaḥ ārogyasya jñānam asti.} — Ayurveda is the knowledge of health.
  \item \emph{Recall:} say the Sanskrit for a decoction, then the Sanskrit for smears, then say \emph{āyurvedaḥ} again
\end{itemize}

\begin{tcolorbox}[breakable,colback=teal!4,colframe=teal!35!black,title={Before you move on}]
What does \sk{आयुर्वेदः} mean? (\textquotedblleft{}The science of health; Ayurveda\textquotedblright{}.) Say it once more.
\end{tcolorbox}
\section[\texorpdfstring{\sk{वातः} (vātaḥ)}{वातः (vātaḥ)}]{\sk{वातः} (vātaḥ) — the wind humour}

\label{lesson:SA-C484-vatah}



\begin{quote}
Before the new one: say the Sanskrit for smears, then the Sanskrit for the science of health.
\end{quote}

\subsection*{You'll want to know: \sk{वातः}}
\textbf{\sk{वातः}} — \emph{vātaḥ} — \textquotedblleft{}the wind humour\textquotedblright{}.

\begin{quote}
\sk{वातः} \sk{शीते} \sk{काले} \sk{वर्धते।}
\end{quote}

\begin{quote}
\emph{vātaḥ śīte kāle vardhate.}
\end{quote}

\begin{quote}
The wind humour rises in the cold season.
\end{quote}

A masculine noun.

\begin{cousinweb}
\textbf{\sk{वातः}} (\emph{vātaḥ}) comes from the root √vā, \textquotedblleft{}to blow\textquotedblright{}; MW gives \textquotedblleft{}wind\textquotedblright{} and \textquotedblleft{}wind or air as one of the humours of the body\textquotedblright{} (MW).
\end{cousinweb}

\begin{grammarlens}[title={What you've built}]
One more word for describing what you lived and felt.
\end{grammarlens}

\subsection*{Your turn}
\begin{itemize}
\raggedright
  \item \emph{Say it:} \emph{vātaḥ}
  \item \emph{Say it:} \emph{vātaḥ}, once more
  \item \emph{Say it:} \emph{vātaḥ śīte kāle vardhate.} — The wind humour rises in the cold season.
  \item \emph{Recall:} say the Sanskrit for smears, then the Sanskrit for the science of health, then say \emph{vātaḥ} again
\end{itemize}

\begin{tcolorbox}[breakable,colback=teal!4,colframe=teal!35!black,title={Before you move on}]
What does \sk{वातः} mean? (\textquotedblleft{}The wind humour\textquotedblright{}.) Say it once more.
\end{tcolorbox}
\section[\texorpdfstring{\sk{पित्तम्} (pittam)}{पित्तम् (pittam)}]{\sk{पित्तम्} (pittam) — bile}

\label{lesson:SA-C484-pittam}



\begin{quote}
Before the new one: say the Sanskrit for the science of health, then the Sanskrit for the wind humour.
\end{quote}

\subsection*{You'll want to know: \sk{पित्तम्}}
\textbf{\sk{पित्तम्}} — \emph{pittam} — \textquotedblleft{}bile\textquotedblright{}.

\begin{quote}
\sk{ग्रीष्मे} \sk{पित्तं} \sk{वर्धते।}
\end{quote}

\begin{quote}
\emph{grīṣme pittaṃ vardhate.}
\end{quote}

\begin{quote}
Bile increases in summer.
\end{quote}

A neuter noun.

\begin{cousinweb}
\textbf{\sk{पित्तम्}} (\emph{pittam}) MW calls its etymology unknown, so its origin is unknown; it gives \textquotedblleft{}bile, the bilious humour, one of the three humours\textquotedblright{} (MW).
\end{cousinweb}

\begin{grammarlens}[title={What you've built}]
One more word for describing what you lived and felt.
\end{grammarlens}

\subsection*{Your turn}
\begin{itemize}
\raggedright
  \item \emph{Say it:} \emph{pittam}
  \item \emph{Say it:} \emph{pittam}, once more
  \item \emph{Say it:} \emph{grīṣme pittaṃ vardhate.} — Bile increases in summer.
  \item \emph{Recall:} say the Sanskrit for the science of health, then the Sanskrit for the wind humour, then say \emph{pittam} again
\end{itemize}

\begin{tcolorbox}[breakable,colback=teal!4,colframe=teal!35!black,title={Before you move on}]
What does \sk{पित्तम्} mean? (\textquotedblleft{}Bile\textquotedblright{}.) Say it once more.
\end{tcolorbox}
\section[\texorpdfstring{\sk{कफः} (kaphaḥ)}{कफः (kaphaḥ)}]{\sk{कफः} (kaphaḥ) — phlegm}

\label{lesson:SA-C484-kaphah}



\begin{quote}
Before the new one: say the Sanskrit for the wind humour, then the Sanskrit for bile.
\end{quote}

\subsection*{You'll want to know: \sk{कफः}}
\textbf{\sk{कफः}} — \emph{kaphaḥ} — \textquotedblleft{}phlegm\textquotedblright{}.

\begin{quote}
\sk{प्रतिश्याये} \sk{कफः} \sk{वर्धते।}
\end{quote}

\begin{quote}
\emph{pratiśyāye kaphaḥ vardhate.}
\end{quote}

\begin{quote}
With a cold, phlegm increases.
\end{quote}

A masculine noun.

\begin{cousinweb}
\textbf{\sk{कफः}} (\emph{kaphaḥ}) MW gives \textquotedblleft{}phlegm, as one of the three humours\textquotedblright{}; MW gives no derivation, so its origin is uncertain (MW).
\end{cousinweb}

\begin{grammarlens}[title={What you've built}]
One more word for describing what you lived and felt.
\end{grammarlens}

\subsection*{Your turn}
\begin{itemize}
\raggedright
  \item \emph{Say it:} \emph{kaphaḥ}
  \item \emph{Say it:} \emph{kaphaḥ}, once more
  \item \emph{Say it:} \emph{pratiśyāye kaphaḥ vardhate.} — With a cold, phlegm increases.
  \item \emph{Recall:} say the Sanskrit for the wind humour, then the Sanskrit for bile, then say \emph{kaphaḥ} again
\end{itemize}

\begin{tcolorbox}[breakable,colback=teal!4,colframe=teal!35!black,title={Before you move on}]
What does \sk{कफः} mean? (\textquotedblleft{}Phlegm\textquotedblright{}.) Say it once more.
\end{tcolorbox}
\section[\texorpdfstring{\sk{पथ्यम्} (pathyam)}{पथ्यम् (pathyam)}]{\sk{पथ्यम्} (pathyam) — a wholesome diet}

\label{lesson:SA-C484-pathyam}



\begin{quote}
Before the new one: say the Sanskrit for bile, then the Sanskrit for phlegm.
\end{quote}

\subsection*{You'll want to know: \sk{पथ्यम्}}
\textbf{\sk{पथ्यम्}} — \emph{pathyam} — \textquotedblleft{}a wholesome diet\textquotedblright{}.

\begin{quote}
\sk{आतुरः} \sk{पथ्यं} \sk{खादेत्।}
\end{quote}

\begin{quote}
\emph{āturaḥ pathyaṃ khādet.}
\end{quote}

\begin{quote}
A patient should eat wholesome food.
\end{quote}

A neuter noun.

\begin{cousinweb}
\textbf{\sk{पथ्यम्}} (\emph{pathyam}) is built on \emph{patha}, \textquotedblleft{}a way\textquotedblright{}: MW gives \textquotedblleft{}belonging to the way\textquotedblright{}, then \textquotedblleft{}suitable, wholesome (esp. of diet)\textquotedblright{} (MW).
\end{cousinweb}

\begin{grammarlens}[title={What you've built}]
One more word for describing what you lived and felt.
\end{grammarlens}

\subsection*{Your turn}
\begin{itemize}
\raggedright
  \item \emph{Say it:} \emph{pathyam}
  \item \emph{Say it:} \emph{pathyam}, once more
  \item \emph{Say it:} \emph{āturaḥ pathyaṃ khādet.} — A patient should eat wholesome food.
  \item \emph{Recall:} say the Sanskrit for bile, then the Sanskrit for phlegm, then say \emph{pathyam} again
\end{itemize}

\begin{tcolorbox}[breakable,colback=teal!4,colframe=teal!35!black,title={Before you move on}]
What does \sk{पथ्यम्} mean? (\textquotedblleft{}A wholesome diet\textquotedblright{}.) Say it once more.
\end{tcolorbox}
